\section{Discussion}

\paragraph{Semantic quality and reliability are different axes.}
Embedding similarity captures whether a completed output remains semantically
close to its source. It does not capture whether the model reliably returns a
usable output. The contrast between Hindi's substantially higher valid-output
coverage and Telugu's conditional scores makes this distinction visible.

\paragraph{Structured generation is part of the task.}
Several failures arise before linguistic quality can be evaluated. Long
responses may exhaust the generation budget before the second JSON field,
while repetition loops consume tokens without completing the requested
structure. Evaluations that silently remove such records measure conditional
quality rather than operational utility.

\paragraph{Automatic metrics are insufficient for naturalness.}
Neither BERTScore nor BGE-M3 evaluates whether a romanization sounds natural to
a Hindi, Telugu, or Kannada speaker. High semantic similarity can coexist with
awkward grammar, inconsistent spelling, inappropriate register, or excessive
English borrowing. Native-speaker judgments are required to establish
practical quality.

\paragraph{Implications for multilingual evaluation.}
Romanized generation should be evaluated with a metric bundle: semantic
preservation, completion coverage, script compliance, repetition rate, and
native-speaker acceptability. Reporting a single mean over successful outputs
can hide the failure mode most visible to users: receiving no usable response.
